% !TEX program = pdflatex
\documentclass[11pt,a4paper]{article}

\usepackage[utf8]{inputenc}
\usepackage[T1]{fontenc}
\usepackage[spanish,es-nodecimaldot]{babel}
\usepackage{amsmath,amssymb,amsthm,mathtools,mathrsfs}
\usepackage{geometry}
\usepackage{enumitem}
\usepackage[expansion=false]{microtype}
\usepackage{hyperref}
\usepackage{cleveref}
\usepackage{xcolor}
\geometry{margin=2.7cm}
\setlength{\parindent}{0pt}
\setlength{\parskip}{0.6em}

\hypersetup{
    colorlinks=true,
    linkcolor=blue!60!black,
    urlcolor=blue!60!black,
    citecolor=blue!60!black
}

%------------------------------------------------
% Entornos
%------------------------------------------------
\theoremstyle{definition}
\newtheorem{definicion}{Definición}[section]
\newtheorem{ejemplo}{Ejemplo}[section]

\theoremstyle{plain}
\newtheorem{proposicion}{Proposición}[section]
\newtheorem{teorema}{Teorema}[section]
\newtheorem{corolario}{Corolario}[section]
\newtheorem{lema}{Lema}[section]

\theoremstyle{remark}
\newtheorem{observacion}{Observación}[section]

%------------------------------------------------
% Comandos
%------------------------------------------------
\newcommand{\C}{\mathbb C}
\newcommand{\R}{\mathbb R}
\newcommand{\Z}{\mathbb Z}
\newcommand{\N}{\mathbb N}
\newcommand{\I}{[0,1]}
\newcommand{\Hol}{\mathcal H}
\newcommand{\OO}{\mathcal O}
\newcommand{\id}{\operatorname{id}}
\newcommand{\restr}[2]{\left.#1\right|_{#2}}

\title{\textbf{Aplicaciones del teorema de monodromía}\\
\large Apuntes sobre la Sección 2.5}
\author{}
\date{}

\begin{document}

\maketitle

\begin{center}
\emph{Basado en la Sección 2.5, ``Applications of the monodromy theorem'',}\\
del Capítulo 2 de \emph{Complex Analysis in One Variable}.
\end{center}

\tableofcontents
\newpage

%================================================
\section{Introducción}
%================================================

La Sección 2.5 desarrolla varias consecuencias del teorema general de
monodromía. La idea central es que un mapa de recubrimiento permite
levantar curvas y, bajo hipótesis topológicas apropiadas, también permite
levantar mapas continuos.

A partir de estas ideas se obtienen, entre otras, las siguientes
consecuencias:

\begin{enumerate}[label=\arabic*.]
    \item un homeomorfismo local con propiedad de levantamiento de curvas
    es un mapa de recubrimiento;
    \item un mapa desde un espacio simplemente conexo puede levantarse
    a un espacio de recubrimiento;
    \item toda función holomorfa sobre un dominio simplemente conexo
    posee una primitiva;
    \item una función holomorfa que no se anula en un dominio simplemente
    conexo posee un logaritmo holomorfo y, por tanto, raíces holomorfas;
    \item el teorema de Poincaré--Volterra proporciona una condición de
    numerabilidad para los espacios de recubrimiento;
    \item las posibles continuaciones analíticas de un germen a lo largo
    de lazos forman un conjunto numerable.
\end{enumerate}

En todo el capítulo se usa la notación
\[
I=[0,1].
\]

%================================================
\section{Levantamiento de curvas y mapas de recubrimiento}
%================================================

\begin{definicion}[Propiedad de levantamiento de curvas]
Sean $X$ y $X'$ variedades y sea
\[
p:X'\longrightarrow X
\]
un homeomorfismo local. Diremos que $p$ tiene la
\textbf{propiedad de levantamiento de curvas} si para todo
$a'\in X'$ y toda curva continua
\[
\gamma:\I\longrightarrow X,
\qquad
\gamma(0)=p(a'),
\]
existe una curva continua
\[
\widetilde\gamma:\I\longrightarrow X'
\]
tal que
\[
p\circ\widetilde\gamma=\gamma,
\qquad
\widetilde\gamma(0)=a'.
\]
\end{definicion}

La Sección 2.1 establece que todo mapa de recubrimiento tiene esta
propiedad. El resultado siguiente demuestra el recíproco.

%------------------------------------------------
\begin{proposicion}
Sean $X$ y $X'$ variedades conexas y sea
\[
p:X'\longrightarrow X
\]
un homeomorfismo local. Supongamos que $p$ tiene la propiedad de
levantamiento de curvas y que $X$ es simplemente conexo. Entonces $p$
es un homeomorfismo global.
\end{proposicion}

\begin{proof}
Fijemos $a'\in X'$ y escribamos
\[
a=p(a').
\]

\textbf{1. $p$ es sobreyectiva.}

Sea $x\in X$. Como $X$ es conexo y una variedad es arcoconexa,
existe una curva
\[
\gamma:\I\longrightarrow X
\]
tal que
\[
\gamma(0)=a,
\qquad
\gamma(1)=x.
\]
Por la propiedad de levantamiento de curvas existe
\[
\widetilde\gamma:\I\longrightarrow X'
\]
con
\[
p\circ\widetilde\gamma=\gamma,
\qquad
\widetilde\gamma(0)=a'.
\]
Evaluando en $1$ obtenemos
\[
p(\widetilde\gamma(1))
=
\gamma(1)
=
x.
\]
Por tanto, $x$ pertenece a la imagen de $p$. Como $x$ era arbitrario,
$p$ es sobreyectiva.

\medskip

\textbf{2. $p$ es inyectiva.}

Sean $a_1,a_2\in X'$ tales que
\[
p(a_1)=p(a_2)=b.
\]
Como $X'$ es conexo, existe una curva
\[
\widetilde\gamma:\I\longrightarrow X'
\]
que une $a_1$ con $a_2$. Consideremos su proyección
\[
\gamma=p\circ\widetilde\gamma.
\]
Entonces $\gamma$ es un lazo de $b$ en $X$.

Como $X$ es simplemente conexo,(por definición) $\gamma$ es homotópico al lazo
constante en $b$. Sea
\[
\{\gamma_u\}_{u\in\I}
\]
tal homotopía, fijando el punto base $b$.

Para cada $u$ levantamos $\gamma_u$ comenzando en $a_1$. Por la
propiedad de levantamiento obtenemos una familia
\[
\{\widetilde\gamma_u\}_{u\in\I}.
\]
El teorema general de monodromía de la sección 2.4 implica que el
extremo final $\widetilde\gamma_u(1)$ no cambia durante la homotopía.

Para el lazo constante, el único levantamiento que comienza en $a_1$
es el lazo constante en $a_1$. En consecuencia,
\[
\widetilde\gamma_1(1)=a_1.
\]
Por otro lado, para $u=0$ el levantamiento es precisamente
$\widetilde\gamma$, de modo que
\[
\widetilde\gamma_0(1)=a_2.
\]
Por la independencia del extremo final,
\[
a_2=a_1.
\]
Por tanto, $p$ es inyectiva.

Finalmente, un homeomorfismo local es una aplicación abierta. Como
$p$ es biyectiva y abierta, su inversa es continua. Por consiguiente,
$p$ es un homeomorfismo global.
\end{proof}

%------------------------------------------------
\begin{teorema}
Sean $X$ y $X'$ variedades y sea
\[
p:X'\longrightarrow X
\]
un homeomorfismo local. Entonces $p$ es una aplicación de recubrimiento
si y sólo si posee la propiedad de levantamiento de curvas.
\end{teorema}

\begin{proof}
La implicación
\[
\text{recubrimiento}\Longrightarrow
\text{propiedad de levantamiento de curvas}
\]
es el resultado de levantamiento de curvas demostrado anteriormente.

Demostremos el recíproco. Sea $a\in X$. Elegimos un entorno simplemente
conexo $U$ de $a$. Como $X$ es una variedad, podemos tomar $U$
homeomorfo a una bola abierta de $\R^n$.

Escribamos
\[
p^{-1}(U)=\bigcup_j U'_j
\]
como unión de sus componentes conexas.

Fijemos una componente $U'_j$ y un punto $a'\in U'_j$. Sea
\[
\gamma:\I\longrightarrow U
\]
una curva con $\gamma(0)=p(a')$. Por hipótesis existe un levantamiento
\[
\widetilde\gamma:\I\longrightarrow X',
\qquad
\widetilde\gamma(0)=a'.
\]
Como $p\circ\widetilde\gamma = \gamma$ está contenido en $U$, tenemos
\[
\widetilde\gamma(\I)\subset p^{-1}(U).
\]
La imagen de una curva es conexa y contiene a $a'$, por lo que debe
estar contenida en la componente $U'_j$. Así, la restricción
\[
p_j:=\restr{p}{U'_j}:U'_j\longrightarrow U
\]
también posee la propiedad de levantamiento de curvas.

La proposición anterior, aplicada a $p_j$, muestra que $p_j$ es un
homeomorfismo de $U'_j$ sobre $U$. Por tanto,
\[
p^{-1}(U)=\coprod_j U'_j
\]
es una unión disjunta de abiertos cada uno de los cuales se proyecta
homeomórficamente sobre $U$. Esto significa exactamente que $U$ está
uniformemente recubierto.

Como $a$ era arbitrario, $p$ es un mapa de recubrimiento.
\end{proof}

%================================================
\section{Levantamiento de mapas}
%================================================

\begin{teorema}
Sean $X$ y $X'$ variedades y sea
\[
p:X'\longrightarrow X
\]
un mapa de recubrimiento. Sea $Y$ una variedad conexa y simplemente
conexa. Fijemos $a'\in X'$ y escribamos
\[
a=p(a').
\]
Supongamos que
\[
f:Y\longrightarrow X
\]
es continua y que existe $y_0\in Y$ tal que
\[
f(y_0)=a.
\]
Entonces existe una única aplicación continua
\[
f':Y\longrightarrow X'
\]
tal que
\[
p\circ f'=f,
\qquad
f'(y_0)=a'.
\]
\end{teorema}

\begin{proof}
Fijemos $y\in Y$. Como $Y$ es conexa, existe una curva
\[
\gamma:\I\longrightarrow Y
\]
que une $y_0$ con $y$:
\[
\gamma(0)=y_0,
\qquad
\gamma(1)=y.
\]
La composición
\[
\alpha=f\circ\gamma
\]
es una curva en $X$ que comienza en $a$.

Como $p$ es un recubrimiento, $\alpha$ posee un levantamiento
\[
\widetilde\alpha:\I\longrightarrow X'
\]
que comienza en $a'$. Definimos
\[
f'(y):=\widetilde\alpha(1).
\]

\medskip
\textbf{Independencia de la curva.}

Supongamos que $\gamma_0$ y $\gamma_1$ son dos curvas de $y_0$ a $y$.
Como $Y$ es simplemente conexo, existe una homotopía
\[
\{\gamma_u\}_{u\in\I}
\]
entre $\gamma_0$ y $\gamma_1$ que fija los extremos.

Al componer con $f$, obtenemos una homotopía
\[
\{f\circ\gamma_u\}_{u\in\I}
\]
entre $f\circ\gamma_0$ y $f\circ\gamma_1$, también fijando los
extremos.

Todos estos caminos parten de $a'$ después de ser levantados.
Por el teorema general de monodromía, sus levantamientos tienen el
mismo extremo final. Por tanto,
\[
\widetilde\alpha_0(1)=\widetilde\alpha_1(1).
\]
Así, $f'(y)$ no depende de la curva elegida.

Por construcción,
\[
f'(y_0)=a'
\]
y
\[
p(f'(y))=f(y).
\]
Queda demostrar la continuidad de $f'$.

\medskip
\textbf{Continuidad.}

Fijemos $y_1\in Y$ y sea $\gamma_1$ una curva de $y_0$ a $y_1$.
Definamos
\[
\alpha_1=f\circ\gamma_1
\]
y sea $\widetilde\alpha_1$ su levantamiento que comienza en $a'$.
Entonces
\[
x_1=f(y_1),
\qquad
x'_1=f'(y_1)=\widetilde\alpha_1(1).
\]

Como $p$ es un homeomorfismo local, existen entornos arcoconexos
$U'$ de $x'_1$ y $U$ de $x_1$ tales que
\[
p|_{U'}:U'\longrightarrow U
\]
es un homeomorfismo.

Por continuidad de $f$, podemos elegir un entorno arcoconexo $V$ de
$y_1$ tal que
\[
f(V)\subset U.
\]

Sea $y\in V$. Elegimos una curva $\delta$ contenida en $V$ que une
$y_1$ con $y$. Entonces la concatenación
\[
\gamma_1*\delta
\]
es una curva de $y_0$ a $y$. La imagen mediante $f$ es
\[
(f\circ\gamma_1)*(f\circ\delta).
\]
Su levantamiento que comienza en $a'$ está dado por
\[
\widetilde\alpha_1*
\left((p|_{U'})^{-1}\circ f\circ\delta\right).
\]
En particular,
\[
f'(y)
=
(p|_{U'})^{-1}(f(y)).
\]
El lado derecho es continuo como composición de aplicaciones
continuas. Por tanto, $f'$ es continua en $y_1$.

Como $y_1$ era arbitrario, $f'$ es continua en todo $Y$.

La unicidad sigue del lema de unicidad de levantamientos: si $f'_1$ y
$f'_2$ son dos levantamientos de $f$ que coinciden en $y_0$, entonces
coinciden en todo $Y$, pues $Y$ es conexo.
\end{proof}

%================================================
\section{Primitivas en dominios simplemente conexos}
%================================================

\begin{teorema}
Sea $\Omega\subset\C$ un abierto simplemente conexo. Entonces toda
función
\[
f\in\Hol(\Omega)
\]
admite una primitiva en $\Omega$.
\end{teorema}

\begin{proof}
La demostración puede obtenerse directamente del teorema de Cauchy en
su forma homotópica. Aquí se presenta la construcción mediante el
teorema de levantamiento de mapas.

Consideremos la aplicación
\[
\varphi:\Omega\longrightarrow\mathcal O,
\]
que a cada $z\in\Omega$ asigna el germen de $(\Omega,f)$ en $z$.
Es decir,
\[
\varphi(z)=\text{germen de }f\text{ en }z.
\]

Sea
\[
p:\mathcal O\longrightarrow\C
\]
la proyección natural que envía un germen al punto donde está definido.
Por la teoría de la gavilla de gérmenes, $p$ es un mapa de recubrimiento
sobre la componente correspondiente.

Como $\Omega$ es simplemente conexo, el teorema anterior permite
levantar $\varphi$ a una aplicación continua
\[
\Phi:\Omega\longrightarrow\mathcal O
\]
tal que
\[
p\circ\Phi=\varphi.
\]

Definimos
\[
F(z):=\Phi(z)(z),
\]
es decir, $F(z)$ es el valor en $z$ del germen $\Phi(z)$.

Demostremos que $F$ es holomorfa. Sea $z_0\in\Omega$ y supongamos que
$(U,g)$ representa al germen $\Phi(z_0)$. Por la definición de la
topología de la gavilla y la continuidad de $\Phi$, existe un entorno
$V$ de $z_0$ tal que, para todo $z\in V$,
\[
\Phi(z)
\]
es el germen en $z$ de la misma función $g$. Por tanto,
\[
F|_V=g|_V.
\]
Así, $F$ es holomorfa.

Finalmente, cada germen $\Phi(z)$ pertenece a la clase de continuación
analítica de $f$, por lo que localmente $F$ coincide con $f$ hasta una
constante. En particular,
\[
F'=f.
\]
Por consiguiente, $F$ es una primitiva de $f$.
\end{proof}

\begin{observacion}
El resultado anterior también puede obtenerse de forma más directa a
partir de la forma homotópica del teorema de Cauchy: si dos caminos con
los mismos extremos son homotópicos con extremos fijos, sus integrales
de $f$ coinciden. En un dominio simplemente conexo, esto permite definir
\[
F(z)=\int_{\gamma}f(\zeta)\,d\zeta
\]
independientemente del camino $\gamma$ que una un punto base fijo con
$z$.
\end{observacion}

%------------------------------------------------
\begin{proposicion}
Sea $\Omega\subset\C$ simplemente conexo y sea
\[
f\in\Hol(\Omega)
\]
una función que no se anula en $\Omega$. Entonces existe
\[
g\in\Hol(\Omega)
\]
tal que
\[
e^{g}=f.
\]
La función $g$ es única salvo adición de una constante de la forma
$2\pi i n$, con $n\in\Z$.
\end{proposicion}

\begin{proof}
Como $f$ no se anula,
\[
\frac{f'}{f}
\]
es holomorfa en $\Omega$. Por el teorema anterior existe una primitiva
$\phi$ tal que
\[
\phi'=\frac{f'}{f}.
\]
Entonces
\[
\left(fe^{-\phi}\right)'
=
f'e^{-\phi}-fe^{-\phi}\phi'
=
e^{-\phi}\left(f'-f\frac{f'}{f}\right)
=
0.
\]
Por tanto,
\[
fe^{-\phi}=c
\]
para alguna constante $c\neq0$. Elegimos $\alpha\in\C$ tal que
\[
e^\alpha=c.
\]
Entonces
\[
f=e^{\phi+\alpha}.
\]
Tomando
\[
g=\phi+\alpha
\]
obtenemos
\[
e^g=f.
\]

Si $g_1$ y $g_2$ son dos soluciones, entonces
\[
e^{g_1-g_2}=1.
\]
Por tanto,
\[
g_1-g_2\in2\pi i\Z.
\]
Como $g_1-g_2$ es continua y $\Omega$ es conexo, dicho entero es
constante. Así,
\[
g_1-g_2=2\pi i n
\]
para algún $n\in\Z$.
\end{proof}

%------------------------------------------------
\begin{proposicion}
Sea $\Omega\subset\C$ un dominio simplemente conexo y sea
\[
f\in\Hol(\Omega)
\]
una función que no se anula. Entonces, para todo
\[
n\in\Z,\qquad n\neq0,
\]
existe $h\in\Hol(\Omega)$ tal que
\[
h^n=f.
\]
\end{proposicion}

\begin{proof}
Por la proposición anterior existe $g\in\Hol(\Omega)$ tal que
\[
e^g=f.
\]
Definimos
\[
h=e^{g/n}.
\]
Entonces
\[
h^n
=
e^g
=
f.
\]
\end{proof}

%================================================
\section{El teorema de Poincaré--Volterra}
%================================================

El siguiente resultado es de naturaleza topológica. Su finalidad en
esta sección es garantizar que ciertos espacios de recubrimiento
relevantes en análisis complejo tienen una base numerable.

\begin{teorema}[Poincaré--Volterra]
Sea $X$ una variedad conexa y sea $Y$ un espacio topológico de Hausdorff
que posee una base numerable. Supongamos que existe una aplicación
continua
\[
\phi:X\longrightarrow Y
\]
tal que, para todo $y\in Y$, la fibra
\[
\phi^{-1}(y)
\]
es discreta en $X$. Entonces $X$ posee una base numerable.
\end{teorema}

\begin{proof}
Sea
\[
\{V_\nu\}_{\nu\in\N}
\]
una base numerable de $Y$.

Consideremos la familia $\mathcal F$ de los abiertos no vacíos
$U\subset X$ que satisfacen:

\begin{enumerate}[label=(\roman*)]
    \item $\overline U$ es compacta en $X$;
    \item existe $\nu\in\N$ tal que $U$ es una componente conexa de
    $\phi^{-1}(V_\nu)$.
\end{enumerate}

Demostraremos dos afirmaciones:

\begin{enumerate}[label=\textbf{Parte \arabic*.}]
    \item $\mathcal F$ es una base de $X$;
    \item $\mathcal F$ es numerable.
\end{enumerate}

\subsubsection*{Parte 1: $\mathcal F$ es una base}

Sea $a\in X$ y sea $W$ un entorno abierto de $a$. Debemos encontrar
$U\in\mathcal F$ tal que
\[
a\in U\subset W.
\]

Como $\phi^{-1}(\phi(a))$ es discreta y $X$ es localmente compacto,
podemos escoger un abierto $W'\subset W$ tal que
\[
a\in W',
\qquad
\overline{W'}\text{ es compacta},
\qquad
\overline{W'}\cap\phi^{-1}(\phi(a))=\{a\}.
\]
Entonces
\[
K:=\phi(\overline{W'}\setminus W')
\]
es compacto y no contiene a $\phi(a)$. Como $\{V_\nu\}$ es una base de
$Y$, podemos elegir $\nu$ tal que
\[
\phi(a)\in V_\nu,
\qquad
V_\nu\cap K=\varnothing.
\]

Sea $U$ la componente conexa de
\[
\phi^{-1}(V_\nu)
\]
que contiene a $a$.

Afirmamos que
\[
U\subset W'.
\]
En efecto, si
\[
x\in U\cap(\overline{W'}\setminus W'),
\]
entonces
\[
\phi(x)\in K.
\]
Pero $x\in\phi^{-1}(V_\nu)$, de modo que
\[
\phi(x)\in V_\nu\cap K,
\]
lo cual contradice la elección de $V_\nu$.

Así,
\[
U\cap(\overline{W'}\setminus W')=\varnothing.
\]
Como $U$ es conexo y contiene a $a\in W'$, no puede atravesar la
frontera de $W'$. Por tanto,
\[
U\subset W'.
\]
Además, $U$ es relativamente compacto porque $W'$ lo es. Por
construcción, $U$ es una componente de $\phi^{-1}(V_\nu)$.
Consecuentemente,
\[
U\in\mathcal F.
\]
Hemos demostrado
\[
a\in U\subset W.
\]
Por tanto, $\mathcal F$ es una base de $X$.

\subsubsection*{Parte 2: $\mathcal F$ es numerable}

Fijemos $U_0\in\mathcal F$ y definamos
\[
\mathcal F_0=\{U_0\}.
\]
Inductivamente, para $k\geq0$, definimos
\[
\mathcal F_{k+1}
=
\left\{
U\in\mathcal F:
\exists\,U'\in\mathcal F_k
\text{ con }
U\cap U'\neq\varnothing
\right\}.
\]

Afirmamos primero que
\[
\mathcal F=\bigcup_{k\geq0}\mathcal F_k.
\]
En efecto, sea
\[
\mathcal F'=\bigcup_{k\geq0}\mathcal F_k
\]
y sea
\[
\mathcal F''=\mathcal F\setminus\mathcal F'.
\]
Los conjuntos
\[
X'=\bigcup_{U\in\mathcal F'}U,
\qquad
X''=\bigcup_{U\in\mathcal F''}U
\]
son abiertos y
\[
X=X'\cup X''.
\]
Además, $X'\cap X''=\varnothing$. Si hubiera un punto en la
intersección, existirían
\[
U'\in\mathcal F',
\qquad
U''\in\mathcal F''
\]
con
\[
U'\cap U''\neq\varnothing.
\]
Pero $U'\in\mathcal F_k$ para algún $k$, lo que implicaría
\[
U''\in\mathcal F_{k+1},
\]
contradicción. Como $X$ es conexo y $X'\neq\varnothing$, obtenemos
$X''=\varnothing$. Por tanto,
\[
\mathcal F=\bigcup_{k\geq0}\mathcal F_k.
\]

Basta demostrar que cada $\mathcal F_k$ es numerable.

Para $k=0$ es evidente.

Supongamos que $\mathcal F_k$ es numerable. Sea
\[
X_k=\bigcup_{U\in\mathcal F_k}U.
\]
Cada elemento de $\mathcal F$ es relativamente compacto y, al ser
una variedad, admite una base numerable. Como $\mathcal F_k$ es
numerable, $X_k$ también posee una base numerable.

Fijemos $\nu\in\N$ y consideremos
\[
\mathcal E_\nu
=
\left\{
U\in\mathcal F_{k+1}:
U\text{ es componente conexa de }\phi^{-1}(V_\nu)
\right\}.
\]
Los conjuntos
\[
U\cap X_k,
\qquad U\in\mathcal E_\nu,
\]
son abiertos no vacíos y dos a dos disjuntos en $X_k$. En un espacio
con base numerable, toda familia de abiertos no vacíos y dos a dos
disjuntos es numerable: a cada abierto se le puede asociar un elemento
distinto de una base numerable contenido en él.

Por tanto, $\mathcal E_\nu$ es numerable para cada $\nu$. Como
\[
\mathcal F_{k+1}
=
\bigcup_{\nu\in\N}\mathcal E_\nu,
\]
se sigue que $\mathcal F_{k+1}$ es numerable.

Por inducción, cada $\mathcal F_k$ es numerable. Como
$\mathcal F$ es unión numerable de familias numerables, $\mathcal F$
es numerable.

Ya vimos que $\mathcal F$ es una base de $X$. Por consiguiente, $X$
tiene una base numerable.
\end{proof}

%================================================
\section{Aplicación: continuaciones analíticas}
%================================================

\begin{corolario}
Sea $a\in\C$ y sea
\[
\varphi_a\in\mathcal O_a
\]
un germen de función holomorfa en $a$. Sea $A$ el conjunto de todos
los elementos de $\mathcal O_a$ que se obtienen mediante continuación
analítica de $\varphi_a$ a lo largo de algún lazo basado en $a$.
Entonces $A$ es numerable.
\end{corolario}

\begin{proof}
Sea
\[
p:\mathcal O\longrightarrow\C
\]
la proyección natural de la gavilla de gérmenes:
\[
p(\varphi_z)=z.
\]
Sea $X$ la componente conexa de $\mathcal O$ que contiene a
$\varphi_a$.

Por la teoría de la gavilla, $p$ es un homeomorfismo local. Además,
las fibras
\[
p^{-1}(z)\cap X
\]
son discretas en $X$: dos gérmenes distintos sobre el mismo punto
pueden separarse mediante entornos de gérmenes disjuntos.

El espacio $\C$ tiene una base numerable. Por el teorema de
Poincaré--Volterra, $X$ posee entonces una base numerable.

Ahora bien, los gérmenes obtenidos por continuación analítica a lo
largo de lazos basados en $a$ son exactamente
\[
A=p^{-1}(a)\cap X.
\]
Esta fibra es discreta en $X$. En un espacio con base numerable, todo
subconjunto discreto es numerable. Por consiguiente,
\[
A\quad\text{es numerable}.
\]
\end{proof}

%================================================
\section{El teorema clásico de monodromía como aplicación}
%================================================

\begin{teorema}[Teorema clásico de monodromía]
Sean
\[
\gamma_0,\gamma_1:\I\longrightarrow\C
\]
dos curvas con los mismos extremos:
\[
\gamma_0(0)=\gamma_1(0)=a,
\qquad
\gamma_0(1)=\gamma_1(1)=b.
\]
Supongamos que $\{\gamma_u\}_{u\in\I}$ es una homotopía entre
$\gamma_0$ y $\gamma_1$ que fija los extremos.

Sea
\[
\varphi_a\in\mathcal O_a.
\]
Supongamos que $\varphi_a$ puede continuarse analíticamente a lo largo
de $\gamma_u$ para todo $u\in\I$. Entonces las continuaciones de
$\varphi_a$ a lo largo de $\gamma_0$ y $\gamma_1$ determinan el mismo
germen en $b$.
\end{teorema}

\begin{proof}
Sea
\[
p:\mathcal O\longrightarrow\C
\]
la proyección natural de la gavilla de gérmenes. Una continuación
analítica de $\varphi_a$ a lo largo de $\gamma_u$ es precisamente un
levantamiento de $\gamma_u$ mediante $p$ que comienza en
$\varphi_a$.

Por hipótesis, todos los caminos de la homotopía admiten tales
levantamientos y todos comienzan en el mismo punto $\varphi_a$.
El teorema general de monodromía de la sección 2.4 afirma entonces que
los levantamientos forman una homotopía.

Como la homotopía original fija el extremo $b$, el extremo final del
levantamiento es constante respecto de $u$. En particular, los
levantamientos de $\gamma_0$ y $\gamma_1$ tienen el mismo extremo.
Ese extremo es precisamente el germen obtenido por continuación
analítica al final del camino.

Por tanto, las dos continuaciones producen el mismo germen en $b$.
\end{proof}

%================================================
\section{Consecuencias geométricas}
%================================================

\begin{observacion}
Las ideas anteriores explican varias propiedades básicas de los
espacios simplemente conexos.

\begin{enumerate}
    \item Si $X$ es homeomorfo a un espacio simplemente conexo, entonces
    $X$ es simplemente conexo.

    \item Si $X$ y $Y$ son simplemente conexos, entonces
    $X\times Y$ es simplemente conexo.

    \item Todo abierto convexo $\Omega\subset\R^n$ es simplemente
    conexo. Si $\gamma_0,\gamma_1$ tienen los mismos extremos, una
    homotopía explícita está dada por
    \[
    F(t,u)=(1-u)\gamma_0(t)+u\gamma_1(t).
    \]
    La convexidad garantiza que
    \[
    F(t,u)\in\Omega.
    \]

    \item Toda variedad posee, alrededor de cada punto, un entorno
    simplemente conexo, pues puede elegirse un entorno homeomorfo a
    una bola abierta de $\R^n$.

    \item El anillo
    \[
    \{z\in\C:0<|z|<1\}
    \]
    no es simplemente conexo. Esta afirmación puede deducirse de las
    aplicaciones del teorema de monodromía desarrolladas en esta
    sección.
\end{enumerate}
\end{observacion}

%================================================
\section{Resumen lógico de la sección}
%================================================

La estructura conceptual de la Sección 2.5 puede resumirse en la cadena
\[
\boxed{
\text{levantamiento de curvas}
\Longrightarrow
\text{recubrimiento}
\Longrightarrow
\text{levantamiento de mapas}
}
\]
cuando se imponen las hipótesis topológicas adecuadas.

En particular:
\[
\boxed{
X\text{ simplemente conexo}
\Longrightarrow
\text{los caminos no dependen de su deformación}
}
\]
y esto permite construir objetos globales a partir de datos locales.

Para funciones holomorfas:
\[
\boxed{
f\in\Hol(\Omega),\quad
\Omega\text{ simplemente conexo}
\Longrightarrow
f\text{ tiene primitiva}.
}
\]

Si además $f$ no se anula:
\[
\boxed{
f\neq0
\Longrightarrow
f=e^g
\Longrightarrow
f\text{ posee raíces holomorfas}.
}
\]

Finalmente, la construcción de la superficie de gérmenes y el teorema
de Poincaré--Volterra proporcionan:
\[
\boxed{
\text{las posibles continuaciones analíticas alrededor de un punto
forman un conjunto numerable}.
}
\]

\end{document}
